\documentclass[preview]{standalone}

\usepackage{amsmath}
\usepackage{amssymb}
\usepackage{stellar}
\usepackage{definitions}
\usepackage{bettelini}

\begin{document}

\id{classical-stokes-theorem}
\genpage

\section{From Green to Stokes}

\begin{snippet}{green-gauss-to-stokes-interpretation}
    Regard \(\Omega\subseteq\realnumbers^2\) as a surface in the plane
    \(z=0\), with normal \(\mathbf k\), and associate to
    \(P\,dx+Q\,dy\) the field \(F=(P,Q,0)\). Then
    \[
        \operatorname{curl}F=(0,0,Q_x-P_y),
    \]
    so Green--Gauss becomes
    \[
        \iint_\Omega\operatorname{curl}F\cdot\mathbf k\,dS
        =\int_{\partial\Omega}F\cdot T\,d\sigma.
    \]
    Stokes' theorem extends this identity to oriented surfaces in
    \(\realnumbers^3\).
\end{snippet}

\begin{snippettheorem}{classical-stokes-theorem-statement}{Stokes' theorem}
    Let \(\Sigma\subseteq\realnumbers^3\) be an oriented, piecewise regular
    surface with piecewise regular boundary \(\partial\Sigma\), and let
    \(F\in\mathcal C^1(U,\realnumbers^3)\) on an open neighborhood of
    \(\Sigma\). If \(\nu\) is the chosen unit normal and the boundary has
    the induced positive orientation, then
    \[
        \iint_\Sigma\operatorname{curl}F\cdot\nu\,dS
        =\int_{\partial\Sigma}F\cdot T\,d\sigma.
    \]
\end{snippettheorem}

\begin{snippet}{stokes-boundary-orientation}
    The boundary orientation is induced by the chosen surface normal.
    Looking from the tip of \(\nu\), a positively oriented boundary is
    traversed counterclockwise. Reversing \(\nu\) reverses the boundary
    orientation and changes both sides of Stokes' formula by a sign.
\end{snippet}

\section{Examples}

\begin{snippetexample}{stokes-spatial-differential-form-example}{A line integral along a spatial curve}
    Let
    \[
        \omega=(y+z)\,dx+(z+x)\,dy+(x-y)\,dz,
    \]
    and let \(\gamma\) be the intersection
    \[
        x^2+y^2+z^2=1,
        \qquad z=y.
    \]
    Compute \(\int_\gamma\omega\).
\end{snippetexample}

\begin{snippetsolution}{stokes-spatial-differential-form-solution}{A line integral along a spatial curve}
    The curve satisfies \(x^2+2y^2=1\) and can be parametrized by
    \[
        \gamma(\theta)=
        \left(\cos\theta,
        \frac{\sin\theta}{\sqrt2},
        \frac{\sin\theta}{\sqrt2}\right),
        \qquad0\leq\theta\leq2\pi.
    \]
    We have
    \[
        dx=-\sin\theta\,d\theta,
        \qquad
        dy=dz=\frac{\cos\theta}{\sqrt2}\,d\theta.
    \]
    Direct substitution gives
    \begin{align*}
        \int_\gamma\omega
        &=\int_0^{2\pi}\Bigg[
            \sqrt2\sin\theta(-\sin\theta)
            +\left(\frac{\sin\theta}{\sqrt2}+\cos\theta\right)
                \frac{\cos\theta}{\sqrt2}
            +\left(\cos\theta-\frac{\sin\theta}{\sqrt2}\right)
                \frac{\cos\theta}{\sqrt2}
        \Bigg],d\theta\\
        &=\sqrt2\integral[0][2\pi]
        [\cos^2\theta-\sin^2\theta][\theta]\\
        &=0.
    \end{align*}

    Alternatively, associate
    \(F=(y+z,z+x,x-y)\) to \(\omega\), and let \(\Sigma\) be the planar
    region in \(z=y\) bounded by \(\gamma\). A unit normal is
    \(\nu=(0,1,-1)/\sqrt2\), while
    \[
        \operatorname{curl}F=(-2,0,0).
    \]
    Hence \(\operatorname{curl}F\cdot\nu=0\), and Stokes' theorem again
    yields \(\int_\gamma\omega=0\).
\end{snippetsolution}

\begin{snippetexample}{stokes-graph-curl-flux-example}{Flux of a curl through a graph}
    Let \(F(x,y,z)=(z,x,y)\). Compute the flux of
    \(\operatorname{curl}F\) through
    \[
        \Sigma=\{(x,y,z):z=xy,\ x^2+y^2\leq1\},
    \]
    oriented by the normal with nonnegative third component.
\end{snippetexample}

\begin{snippetsolution}{stokes-graph-curl-flux-solution}{Flux of a curl through a graph}
    The parametrization \(r(x,y)=(x,y,xy)\) induces the normal
    \(r_x\times r_y=(-y,-x,1)\). Its positively oriented boundary is
    \[
        \gamma(\theta)=
        (\cos\theta,\sin\theta,\cos\theta\sin\theta),
        \qquad0\leq\theta\leq2\pi.
    \]
    Thus
    \[
        \gamma'(\theta)=
        (-\sin\theta,\cos\theta,\cos^2\theta-\sin^2\theta)
    \]
    and
    \(F(\gamma(\theta))=(\cos\theta\sin\theta,\cos\theta,\sin\theta)\).
    Stokes' theorem gives
    \begin{align*}
        \iint_\Sigma\operatorname{curl}F\cdot\nu\,dS
        &=\integral[0][2\pi]
        [F(\gamma(\theta))\cdot\gamma'(\theta)][\theta]\\
        &=\integral[0][2\pi]
        [-\sin^2\theta\cos\theta
          +\cos^2\theta
          +\sin\theta(\cos^2\theta-\sin^2\theta)][\theta]\\
        &=\integral[0][2\pi][\cos^2\theta][\theta]=\pi,
    \end{align*}
    because all remaining terms integrate to zero over a full period.
\end{snippetsolution}

\end{document}
